\documentclass[9pt,a4paper]{extarticle}
\usepackage[T1]{fontenc}
\usepackage[utf8]{inputenc}
\usepackage[margin=1.25cm]{geometry}
\usepackage{amsmath,booktabs,tabularx,array,makecell}
\usepackage[table]{xcolor}
\usepackage{pgfplotstable}
\pgfplotsset{compat=1.18}
\usepackage[protrusion=true,expansion=false]{microtype}
\usepackage{enumitem}
\usepackage{fancyhdr}
\setlength{\parindent}{0pt}
\setlength{\parskip}{3pt}
\setlist[itemize]{leftmargin=1.35em,itemsep=1pt,topsep=2pt}
\pagestyle{fancy}
\fancyhf{}
\fancyfoot[C]{\thepage}
\renewcommand{\headrulewidth}{0pt}

% Numerical cells are read directly from the validated CSV artifacts.
\pgfplotstableread[col sep=comma]{comparison_results.csv}\ResultTable
\pgfplotstableread[col sep=comma]{parameter_summary.csv}\ParameterTable
\pgfplotstableread[col sep=comma]{complexity_flops.csv}\FlopTable
\pgfplotstableread[col sep=comma]{selected_hyperparameters.csv}\HyperTable
\newcommand{\Rdec}[2]{\pgfplotstablegetelem{#1}{#2}\of\ResultTable%
  \pgfmathprintnumber[fixed,precision=3]{\pgfplotsretval}}
\newcommand{\Pint}[2]{\pgfplotstablegetelem{#1}{#2}\of\ParameterTable%
  \pgfmathprintnumber[fixed,precision=0]{\pgfplotsretval}}
\newcommand{\Fint}[2]{\pgfplotstablegetelem{#1}{#2}\of\FlopTable%
  \pgfmathprintnumber[fixed,precision=0]{\pgfplotsretval}}
\newcommand{\Hint}[2]{\pgfplotstablegetelem{#1}{#2}\of\HyperTable%
  \pgfmathprintnumber[fixed,precision=0]{\pgfplotsretval}}
\newcommand{\mainrow}{\rowcolor{blue!8}\bfseries}

\begin{document}
\begin{center}
  {\Large\bfseries Fair PN-IQ, GMP and PNNN Comparison}\par
  \vspace{2pt}
  {\normalsize Validated run: \texttt{20260714\_013938}}
\end{center}

\section*{Objective and methodology}
The study compares a complex generalized memory polynomial (GMP), its
phase-normalized real I/Q formulations, and phase-normalized neural networks
(PNNNs) under one frozen data partition and one temporal complex-NMSE
definition.  Coupled PN-IQ is an exact real representation of the selected
complex GMP; Independent PN-IQ relaxes that I/Q coupling; the sparse PNNN is
trained as a wider N12 network and magnitude-pruned to the Independent PN-IQ
active-parameter target.

\section*{Fair comparison protocol}
All models use the same 491,520-sample capture and the local modeled-block
X/Y convention (\texttt{mappingMode=xy\_forward}, without assigning an
additional PA-forward meaning).  DC is removed once.  The deterministic
\texttt{FIT\_POOL} contains 19,661 samples (4.00004\%); its disjoint
\texttt{TEST} complement contains 471,859.  \texttt{TRAIN} (16,711) and
\texttt{VAL} (2,950) partition \texttt{FIT\_POOL}.  Support and lambda are
selected from TRAIN/VAL; linear supports and coefficients are regenerated on
FIT\_POOL.  PNNN normalization and gradient updates use TRAIN, and VAL alone
selects dense and fine-tuned checkpoints; PNNN is not refitted on FIT\_POOL.
TEST is used once after every final model is frozen.  The selected linear
lambdas are \Hint{0}{SelectedLambda} in this run.

\section*{Main results}
\setlength{\tabcolsep}{3pt}
\renewcommand{\arraystretch}{1.12}
\begin{tabularx}{\textwidth}{>{\raggedright\arraybackslash}Xrrrrr}
\toprule
Model & \makecell{Active real\\parameters} &
\makecell{Test NMSE\\(dB)} &
\makecell{Dense core\\FLOPs/sample} &
\makecell{Ideal sparse\\FLOPs/sample} &
\makecell{Special ops\\/sample} \\
\midrule
Complex GMP DOMP-100 & \Pint{0}{NumRealParameters} & \Rdec{0}{TestNMSEdB} & \Fint{0}{DenseExecutionCoreFLOPsPerSample} & \Fint{0}{IdealSparseCoreFLOPsPerSample} & \Fint{0}{NumSpecialOpsPerSample} \\
Coupled PN-IQ DOMP-100 & \Pint{1}{NumRealParameters} & \Rdec{1}{TestNMSEdB} & \Fint{1}{DenseExecutionCoreFLOPsPerSample} & \Fint{1}{IdealSparseCoreFLOPsPerSample} & \Fint{1}{NumSpecialOpsPerSample} \\
\mainrow Independent PN-IQ & \Pint{2}{NumRealParameters} & \Rdec{2}{TestNMSEdB} & \Fint{2}{DenseExecutionCoreFLOPsPerSample} & \Fint{2}{IdealSparseCoreFLOPsPerSample} & \Fint{2}{NumSpecialOpsPerSample} \\
\mainrow Complex GMP parameter-matched & \Pint{3}{NumRealParameters} & \Rdec{3}{TestNMSEdB} & \Fint{3}{DenseExecutionCoreFLOPsPerSample} & \Fint{3}{IdealSparseCoreFLOPsPerSample} & \Fint{3}{NumSpecialOpsPerSample} \\
PNNN H4 dense & \Pint{6}{NumRealParameters} & \Rdec{6}{TestNMSEdB} & \Fint{6}{DenseExecutionCoreFLOPsPerSample} & \Fint{6}{IdealSparseCoreFLOPsPerSample} & \Fint{6}{NumSpecialOpsPerSample} \\
PNNN N12 dense & \Pint{7}{NumRealParameters} & \Rdec{7}{TestNMSEdB} & \Fint{7}{DenseExecutionCoreFLOPsPerSample} & \Fint{7}{IdealSparseCoreFLOPsPerSample} & \Fint{7}{NumSpecialOpsPerSample} \\
\mainrow PNNN N12 sparse & \Pint{8}{NumRealParameters} & \Rdec{8}{TestNMSEdB} & \Fint{8}{DenseExecutionCoreFLOPsPerSample} & \Fint{8}{IdealSparseCoreFLOPsPerSample} & \Fint{8}{NumSpecialOpsPerSample} \\
\bottomrule
\end{tabularx}

\vfill
\newpage

\section*{Computational cost}
FLOPs refer to inference per modeled output sample; DOMP is used only during
identification and is therefore excluded.  The core convention assigns one
FLOP to a real addition or multiplication, two to a complex addition, and six
to a complex multiplication.  ELU, worst-case exponential, square-root,
division, and absolute-value calls are reported separately because no
platform-independent conversion was imposed.

For the sparse N12 PNNN, the two execution views are:
\begin{center}
\begin{tabular}{lrrrrr}
\toprule
Execution view & Real mult. & Real add. & Core FLOPs/sample & ELU/sample & Special ops/sample \\
\midrule
Dense matrices & \Fint{8}{RealMultiplicationsPerSample} & \Fint{8}{RealAdditionsPerSample} & \Fint{8}{DenseExecutionCoreFLOPsPerSample} & \Fint{8}{NumELUPerSample} & \Fint{8}{NumSpecialOpsPerSample} \\
Ideal sparse kernel & \Fint{8}{IdealSparseRealMultiplicationsPerSample} & \Fint{8}{IdealSparseRealAdditionsPerSample} & \Fint{8}{IdealSparseCoreFLOPsPerSample} & \Fint{8}{NumELUPerSample} & \Fint{8}{NumSpecialOpsPerSample} \\
\bottomrule
\end{tabular}
\end{center}
The dense count executes the full N12 matrices, including stored zeros.  The
ideal sparse count includes only active weights and is attainable only if the
implementation actually skips zero products; pruning alone does not guarantee
this behavior in MATLAB or hardware.  The 12 ELUs remain in both views.

\section*{Short interpretation}
\begin{itemize}
  \item Independent PN-IQ gives the best TEST NMSE among the highlighted
  parameter-matched models: $-38.492$~dB with 358 active real parameters.
  The closest complex GMP uses 352 parameters and reaches $-38.126$~dB;
  sparse N12 uses exactly 358 and reaches $-35.457$~dB.
  \item Sparse N12 has a lower \emph{ideal} core count (876) than Independent
  PN-IQ (1062), but a higher dense-execution count (2252).  Its nonlinear
  activations also add unweighted special operations.
  \item The coupled and complex DOMP-100 predictions agree to a relative error
  of $1.320\times10^{-11}$, confirming a representation change rather than an
  NMSE improvement.
  \item No single model wins every metric: PN-IQ leads NMSE here, whereas an
  implementation-aware sparse PNNN could lead ideal arithmetic count.
\end{itemize}

\section*{Limitations and traceability}
Results are from one capture, one neural seed, a 4\% fit pool, and analytical
operation counts.  They do not establish universal accuracy, physical
latency, energy, FPGA resources, or hardware superiority.  The experiment
does not compare quantization, synthesis, or valid demodulated ACPR/EVM.
Every number above is read from the validated
\texttt{comparison\_results.csv}, \texttt{parameter\_summary.csv},
\texttt{complexity\_flops.csv}, and
\texttt{selected\_hyperparameters.csv}; no experimental model was rerun.

\end{document}
